Una señal normalizada satisface el modelo suministrado
\[
x(t)=e^{-kt},\qquad x(0)=1,\qquad k\ge0,
\]
con \(t\) en segundos. Cada lectura de \(x\) tiene un error absoluto garantizado no mayor que \(0.02\); los tiempos y el valor inicial se consideran exactos. Tres pruebas piloto simuladas del mismo régimen han dado
\[
\begin{array}{c|ccc}
t\ (\mathrm s)&0.10&1.00&4.00\\\hline
x_{\rm obs}&0.90&0.37&0.02
\end{array}
\]
Se quiere elegir uno de esos instantes para la siguiente medición. Considere primero cada lectura por separado, sin combinar las tres, y determine todos los valores de \(k\) compatibles con ella y con su cota de error. Use esos intervalos para elegir el instante que ha proporcionado la acotación más ajustada. Explique la diferencia entre medir cuando la señal apenas ha cambiado y medir cuando se aproxima al nivel del error del instrumento; no se pide encontrar un instante óptimo entre todos los tiempos posibles.
